
\section*{Kurzreferat}

Die Arbeit untersucht, ob ungelabelte Bilder das Deep Metric Learning auf eingefrorenen Merkmalen des Vision Transformers DINOv2 verbessern. Vier semi-supervised Mechanismen (Label Propagation, Laplace-Regularisierung, Entropieminimierung nach Hoffer und Ailon, SLADE) werden mit fünf Verlustfunktionen kombiniert. Jede Kombination wird auf Cars196, CUB, In-Shop und Semi-iNat-2021 mit klassendisjunkter Kreuzvalidierung einzeln optimiert. Bei mehreren Label-Budgets wird sie mit fünf Seeds wiederholt und mit einer gleich optimierten supervised Referenz verglichen. Auf Cars196 und In-Shop verbessert Label Propagation bei den Budgets der Hyperparametersuche die MAP@R eines Einzelmodells in allen gepaarten Vergleichen, auf Cars196 mit 16 gelabelten Bildern je Klasse um 2,2 und auf In-Shop mit zwei um 3,2 Prozentpunkte. Hoffer gewinnt auf In-Shop ähnlich viel. Mit derselben Konfiguration gewinnt Label Propagation auf Cars196 bei vier gelabelten Bildern je Klasse 7,8 Prozentpunkte und erreicht damit die Referenz mit 16. Mit steigendem Budget wird der Gewinn klein oder verschwindet. Auf CUB sind die unveränderten DINOv2-Merkmale besser als die meisten trainierten Einzelmodelle. Kein Verfahren gewinnt dort mehr als einen Prozentpunkt. Auf Semi-iNat-2021 gewinnt die Referenz von zwei bis hundert Labels je Art nur rund 3 Prozentpunkte und nur Label Propagation verbessert sie verlässlich, um 0,6 Prozentpunkte bei zwei Labels je Art. Ablationen auf Cars196 und In-Shop wurden mit einem Seed und allen fünf Verlustfunktionen durchgeführt. Die separat gewählten Hyperparameter allein erklären den Gewinn nicht. Mit den Hyperparametern der Referenz helfen die ungelabelten Bilder vor allem bei Label Propagation. Je nach Verfahren und Verlustfunktion hilft der semi-supervised Term schon ohne ungelabelte Bilder, etwa die Laplace-Regularisierung auf Cars196. Die graphbasierten Verfahren brauchen die aus den Daten berechnete Nachbarschaft. Sie deuten zudem an, dass Fehler der Label Propagation weniger schaden als zufällige Labelfehler gleicher mittlerer Häufigkeit, auf Cars196 mit drei der fünf Verlustfunktionen und auf In-Shop mit allen fünf. Die Verfolgung der Embedding-Geometrie auf Cars196 mit vier Labels je Klasse zeigt, dass die supervised Referenz nur die gelabelten Bilder zusammenzieht. Label Propagation macht nach dem Warm-up auch die ungelabelten Bilder kompakter und ordnet sie häufiger als die übrigen Mechanismen der richtigen Klasse zu.

\clearpage
\section*{Abstract}

This thesis examines whether unlabeled images improve deep metric learning on frozen features of the vision transformer DINOv2. Four semi-supervised mechanisms (label propagation, Laplacian regularization, entropy minimization following Hoffer and Ailon, and SLADE) are combined with five loss functions. Each combination is tuned individually on Cars196, CUB, In-Shop and Semi-iNat-2021 with class-disjoint cross-validation. It is repeated with five seeds at several label budgets and compared with an equally tuned supervised baseline. On Cars196 and In-Shop, label propagation improves the MAP@R of a single model in every paired comparison at the budgets used for the hyperparameter search, by 2.2 percentage points on Cars196 with 16 labeled images per class and by 3.2 points on In-Shop with two. Hoffer's method gains about as much on In-Shop. With the same configuration, label propagation gains 7.8 points on Cars196 with four labeled images per class and matches the baseline trained with 16. The gain shrinks or vanishes as the label budget grows. On CUB, the unchanged DINOv2 features outperform most trained single models, and no method gains more than one point. On Semi-iNat-2021 the baseline gains only about 3 points from two to one hundred labels per species, and only label propagation improves on it consistently, by 0.6 points with two labels per species. Single-seed ablations were performed on Cars196 and In-Shop with all five losses. The separately tuned hyperparameters alone do not explain the gain. With the baseline's hyperparameters, the unlabeled images help mainly for label propagation. Depending on the method and loss, the semi-supervised term already helps without unlabeled images, for example Laplacian regularization on Cars196. The graph-based methods depend on the neighborhood structure computed from the data. They also suggest that label propagation errors hurt less than random label noise of the same average rate, on Cars196 with three of the five losses and on In-Shop with all five. Tracking the embedding geometry on Cars196 with four labels per class shows that the supervised baseline only pulls the labeled images together. After its warm-up, label propagation also makes the unlabeled images more compact and assigns more of them to the correct class than the other mechanisms.
